% Source: fourier_course.tex
\begin{notechapter}
% ---- chapter-local definitions ----
\providecommand{\sinc}{\operatorname{sinc}}
\providecommand{\rect}{\operatorname{rect}}
\providecommand{\fourier}{\mathcal{F}}
% -----------------------------------
\chapter{Fourier Analysis}\label{ch:fourier-analysis}
\chaptersubtitle{Series, transforms and the FFT}

\begin{chapterintro}
Spectral graph learning (Chapter~\ref{ch:graph-laplacians} and the next chapter) borrows its vocabulary from classical Fourier analysis: \emph{frequency}, \emph{filter}, \emph{convolution theorem}, \emph{low-pass}. This chapter supplies that vocabulary. We treat the four members of the Fourier family (Fourier series, the Fourier transform, the discrete-time Fourier transform and the discrete Fourier transform) and the fast algorithm that makes the last one practical. For each we state when it exists, derive its inversion formula, and prove the properties that are used later: Parseval's identity, the shift rules and, above all, the convolution theorem, which says that a transform \emph{diagonalises} convolution. You will learn
\begin{itemize}[nosep]
  \item why the coefficients of a Fourier series are inner products with orthogonal exponentials, and in what sense the series converges;
  \item the Fourier transform as the limit of the series, with Riemann--Lebesgue, inversion and Plancherel;
  \item the shift, scaling, differentiation and convolution rules;
  \item sampling, aliasing and the Nyquist--Shannon theorem through the DTFT;
  \item the DFT as a matrix, circular convolution, and the radix-2 FFT with its $O(N\log N)$ cost.
\end{itemize}
\prerequisites{Lebesgue integral and $L^p$ spaces at the level of definitions; complex exponentials; linear algebra (unitary matrices).}
\end{chapterintro}

\section{The Fourier family}

Every member of the family represents a signal as a superposition of complex exponentials $e^{i\omega t}$. Which member applies depends on two properties of the signal: whether time is continuous or discrete, and whether the signal is periodic. A rule worth remembering is that \emph{periodicity in one domain corresponds to discreteness in the other}.

\begin{center}
\begin{tabular}{lll}
\toprule
\textbf{Transform} & \textbf{Time domain} & \textbf{Frequency domain} \\
\midrule
Fourier series (FS) & continuous, periodic & discrete \\
Fourier transform (FT) & continuous, aperiodic & continuous \\
Discrete-time FT (DTFT) & discrete, aperiodic & continuous, $2\pi$-periodic \\
Discrete FT (DFT) & discrete, $N$-periodic & discrete, $N$-periodic \\
\bottomrule
\end{tabular}
\end{center}

\begin{keyidea}
All four transforms expand a signal in an orthogonal family of complex exponentials. Complex exponentials are the eigenfunctions of every shift-invariant linear operator, which is why each transform turns convolution into pointwise multiplication. On a graph, shifts are replaced by the Laplacian, and its eigenvectors play the role of the exponentials.
\end{keyidea}

\section{Mathematical Preliminaries}

\subsection{Function Spaces}

\begin{definition}[$L^1$ Space]
The space $L^1(\R)$ consists of all Lebesgue measurable functions $f:\R \to \C$ such that
\[
\norm{f}_1 = \int_{-\infty}^{\infty} \abs{f(t)} \, dt < \infty.
\]
\end{definition}

\begin{definition}[$L^2$ Space]
The space $L^2(\R)$ consists of all Lebesgue measurable functions $f:\R \to \C$ such that
\[
\norm{f}_2 = \sqrt{\int_{-\infty}^{\infty} \abs{f(t)}^2 \, dt} < \infty.
\]
This is a Hilbert space with inner product
\[
\inner{f}{g} = \int_{-\infty}^{\infty} f(t) \overline{g(t)} \, dt.
\]
\end{definition}

\begin{theorem}[Hölder's Inequality]
If $p, q \geq 1$ satisfy $\frac{1}{p} + \frac{1}{q} = 1$, then for $f \in L^p$ and $g \in L^q$,
\[
\int_{-\infty}^{\infty} \abs{f(t)g(t)} \, dt \leq \norm{f}_p \norm{g}_q.
\]
\end{theorem}

\begin{pitfall}[$L^1$ versus $L^2$]
On a bounded interval, Cauchy--Schwarz (Hölder with $p=q=2$) gives $\int_0^T \abs{f} \le \sqrt{T}\,\norm{f}_{L^2([0,T])}$, so $L^2([0,T]) \subset L^1([0,T])$. On the whole line neither space contains the other: $f(t) = 1/(1+\abs{t})$ is in $L^2(\R)$ but not in $L^1(\R)$, while $f(t) = \abs{t}^{-1/2}\bbone[\abs{t}<1]$ is in $L^1(\R)$ but not in $L^2(\R)$. This is why the Fourier transform is defined first on $L^1$ and then extended to $L^2$ (Theorem~\ref{fourier-analysis:thm:plancherel}).
\end{pitfall}

\subsection{Periodic Functions}

\begin{definition}[Periodic Function]
A function $f:\R \to \C$ is \emph{$T$-periodic} if $f(t+T) = f(t)$ for all $t \in \R$.
\end{definition}

For periodic functions, we define:
\[
L^2([0,T]) = \left\{ f:[0,T] \to \C : \int_0^T \abs{f(t)}^2 \, dt < \infty \right\}.
\]

\section{Fourier Series}

\subsection{Motivation and Setting}

Consider a $T$-periodic function $f:\R \to \C$. Since the function repeats every $T$ seconds, it is natural to ask whether we can represent it as a sum of harmonically related sinusoids.

\begin{definition}[Fourier Series]
Let $f$ be a $T$-periodic function in $L^2([0,T])$. The \emph{Fourier series} of $f$ is
\[
f(t) \sim \sum_{n=-\infty}^{\infty} c_n e^{i 2\pi n t/T},
\]
where the \emph{Fourier coefficients} are given by
\[
c_n = \frac{1}{T} \int_0^T f(t) e^{-i 2\pi n t/T} \, dt.
\]
\end{definition}

\begin{remark}
We use the symbol $\sim$ to denote that the series is associated with $f$. Convergence requires further conditions.
\end{remark}

\subsection{Orthogonality of Complex Exponentials}

The key insight is that the complex exponentials form an orthogonal basis.

\begin{lemma}[Orthogonality]
For $m, n \in \Z$,
\[
\frac{1}{T}\int_0^T e^{i 2\pi n t/T} e^{-i 2\pi m t/T} \, dt = \delta_{mn} = \begin{cases}
1 & \text{if } m = n \\
0 & \text{if } m \neq n.
\end{cases}
\]
\end{lemma}

\begin{proof}
\[
\frac{1}{T}\int_0^T e^{i 2\pi (n-m) t/T} \, dt = \frac{1}{T} \left[ \frac{T}{i 2\pi (n-m)} e^{i 2\pi (n-m) t/T} \right]_0^T
\]
If $n \neq m$:
\[
= \frac{1}{i 2\pi (n-m)} \left( e^{i 2\pi (n-m)} - 1 \right) = \frac{1}{i 2\pi (n-m)} (1 - 1) = 0.
\]
If $n = m$:
\[
\frac{1}{T}\int_0^T e^{0} \, dt = \frac{1}{T} \cdot T = 1.
\]
\end{proof}

\subsection{Derivation of Fourier Coefficients}

To find the coefficients $c_n$, we use the orthogonality relation. Suppose that $f$ equals its series and that the series may be integrated term by term (true, for instance, when $\sum_n \abs{c_n} < \infty$, by uniform convergence). Multiply both sides of the Fourier series by $e^{-i 2\pi m t/T}$ and integrate:

\begin{align*}
\int_0^T f(t) e^{-i 2\pi m t/T} \, dt &= \int_0^T \left( \sum_{n=-\infty}^{\infty} c_n e^{i 2\pi n t/T} \right) e^{-i 2\pi m t/T} \, dt \\
&= \sum_{n=-\infty}^{\infty} c_n \int_0^T e^{i 2\pi (n-m) t/T} \, dt \\
&= \sum_{n=-\infty}^{\infty} c_n \cdot T \delta_{mn} \\
&= T c_m.
\end{align*}

Therefore,
\[
c_m = \frac{1}{T} \int_0^T f(t) e^{-i 2\pi m t/T} \, dt.
\]

\subsection{Convergence Theorems}

\begin{theorem}[Dirichlet Conditions]
Let $f$ be a $T$-periodic function. If $f$ satisfies:
\begin{enumerate}
\item $f$ is absolutely integrable over one period: $\int_0^T \abs{f(t)} \, dt < \infty$,
\item $f$ has a finite number of maxima and minima in one period,
\item $f$ has a finite number of discontinuities in one period,
\end{enumerate}
then the Fourier series of $f$ converges pointwise to
\[
\frac{f(t^+) + f(t^-)}{2}
\]
at every point $t$, where $f(t^+)$ and $f(t^-)$ denote the right and left limits.
\end{theorem}

\begin{theorem}[Mean Square Convergence]
If $f \in L^2([0,T])$, then
\[
\lim_{N \to \infty} \int_0^T \left| f(t) - \sum_{n=-N}^{N} c_n e^{i 2\pi n t/T} \right|^2 dt = 0.
\]
\end{theorem}

\begin{theorem}[Parseval's Identity for Fourier Series]
For $f \in L^2([0,T])$,
\[
\frac{1}{T} \int_0^T \abs{f(t)}^2 \, dt = \sum_{n=-\infty}^{\infty} \abs{c_n}^2.
\]
\end{theorem}

\begin{proof}
We give the computation when $f$ equals the sum of an absolutely convergent Fourier series, so that sum and integral may be exchanged; the general case follows from Bessel's inequality and the completeness of the exponentials in $L^2([0,T])$, which is the content of the mean-square convergence theorem.
\begin{align*}
\frac{1}{T} \int_0^T \abs{f(t)}^2 \, dt &= \frac{1}{T} \int_0^T f(t) \overline{f(t)} \, dt \\
&= \frac{1}{T} \int_0^T f(t) \overline{\sum_{m} c_m e^{i 2\pi m t/T}} \, dt \\
&= \frac{1}{T} \int_0^T f(t) \sum_{m} \overline{c_m} e^{-i 2\pi m t/T} \, dt \\
&= \sum_{m} \overline{c_m} \frac{1}{T} \int_0^T f(t) e^{-i 2\pi m t/T} \, dt \\
&= \sum_{m} \overline{c_m} c_m = \sum_{m} \abs{c_m}^2.
\end{align*}
\end{proof}

\subsection{Real Fourier Series}

For real-valued functions, we often use the trigonometric form.

\begin{proposition}
If $f:\R \to \R$ is real and $T$-periodic, then $c_{-n} = \overline{c_n}$, and
\[
f(t) = a_0 + \sum_{n=1}^{\infty} \left( a_n \cos\frac{2\pi n t}{T} + b_n \sin\frac{2\pi n t}{T} \right),
\]
where
\begin{align*}
a_0 &= c_0 = \frac{1}{T} \int_0^T f(t) \, dt, \\
a_n &= 2\text{Re}(c_n) = \frac{2}{T} \int_0^T f(t) \cos\frac{2\pi n t}{T} \, dt, \\
b_n &= -2\text{Im}(c_n) = \frac{2}{T} \int_0^T f(t) \sin\frac{2\pi n t}{T} \, dt.
\end{align*}
\end{proposition}

\begin{proof}
For real $f$, conjugating the coefficient formula gives $\overline{c_n} = \frac1T\int_0^T f(t)e^{i2\pi nt/T}dt = c_{-n}$. Write $c_n = (a_n - i b_n)/2$, which is the definition of $a_n$ and $b_n$ above; then $c_{-n} = (a_n + ib_n)/2$. Pairing the terms $n$ and $-n$ and writing $\theta = 2\pi n t/T$,
\begin{align*}
c_n e^{i\theta} + c_{-n} e^{-i\theta}
&= \tfrac12(a_n - ib_n)e^{i\theta} + \tfrac12(a_n + ib_n)e^{-i\theta} \\
&= \tfrac{a_n}{2}\left(e^{i\theta} + e^{-i\theta}\right) - \tfrac{ib_n}{2}\left(e^{i\theta} - e^{-i\theta}\right) \\
&= a_n \cos\theta + b_n \sin\theta,
\end{align*}
using $e^{i\theta}+e^{-i\theta} = 2\cos\theta$ and $e^{i\theta}-e^{-i\theta} = 2i\sin\theta$. The integral formulas for $a_n$ and $b_n$ follow by taking real and imaginary parts of $2c_n$.
\end{proof}

\subsection{Example: Square Wave}

\begin{example}
Consider the square wave with period $T = 2\pi$:
\[
f(t) = \begin{cases}
1 & \text{if } 0 < t < \pi \\
-1 & \text{if } \pi < t < 2\pi
\end{cases}
\]
extended periodically.

The Fourier coefficients are:
\begin{align*}
c_n &= \frac{1}{2\pi} \int_0^{2\pi} f(t) e^{-int} \, dt \\
&= \frac{1}{2\pi} \left( \int_0^{\pi} e^{-int} \, dt - \int_{\pi}^{2\pi} e^{-int} \, dt \right) \\
&= \frac{1}{2\pi} \left[ \frac{e^{-int}}{-in} \right]_0^{\pi} - \frac{1}{2\pi} \left[ \frac{e^{-int}}{-in} \right]_{\pi}^{2\pi} \\
&= \frac{1}{-i 2\pi n} \left( e^{-in\pi} - 1 - e^{-i 2\pi n} + e^{-in\pi} \right) \\
&= \frac{1}{-i 2\pi n} \left( 2e^{-in\pi} - 1 - 1 \right) \\
&= \frac{1}{-i 2\pi n} \left( 2\cos(n\pi) - 2i\sin(n\pi) - 2 \right) \\
&= \frac{1}{-i 2\pi n} \left( 2(-1)^n - 2 \right) \\
&= \frac{(-1)^n - 1}{-i\pi n}.
\end{align*}

For $n$ even (including $n=0$, where the two half-periods cancel): $c_n = 0$.
For $n$ odd: $c_n = \frac{-2}{-i\pi n} = \frac{2}{i\pi n} = -\frac{2i}{\pi n}$.

In the real form, $a_n = 2\operatorname{Re} c_n = 0$ (as it must be, since $f$ is odd) and $b_n = -2\operatorname{Im} c_n = \frac{4}{\pi n}$ for odd $n$. Thus,
\[
f(t) = \sum_{n \text{ odd}} \left(-\frac{2i}{\pi n}\right) e^{int} = \frac{4}{\pi} \sum_{k=0}^{\infty} \frac{\sin((2k+1)t)}{2k+1}.
\]
At the jumps $t = 0, \pi$ every term vanishes, and the series converges to $0$, the average of the one-sided limits $\pm 1$, as the Dirichlet theorem predicts. Since $\abs{c_n} = 2/(\pi\abs{n})$, Parseval gives $\sum_{n \text{ odd}} \frac{4}{\pi^2 n^2} = \frac{1}{2\pi}\int_0^{2\pi} 1\,dt = 1$, that is, $\sum_{k\ge0} (2k+1)^{-2} = \pi^2/8$.
\end{example}

\section{Fourier Transform}

\subsection{Motivation: From Periodic to Aperiodic}

The Fourier transform extends the idea of Fourier series to aperiodic functions by taking the limit as the period $T \to \infty$.

Consider a function $f \in L^1(\R)$. We can create a periodic function $f_T$ that equals $f$ on $[-T/2, T/2]$ and is extended periodically. As $T \to \infty$, $f_T \to f$.

The Fourier series of $f_T$ is:
\[
f_T(t) = \sum_{n=-\infty}^{\infty} c_n e^{i 2\pi n t/T}, \quad c_n = \frac{1}{T} \int_{-T/2}^{T/2} f_T(t) e^{-i 2\pi n t/T} \, dt.
\]

Define $\omega_n = 2\pi n/T$ and $\Delta \omega = 2\pi/T$. Then:
\[
c_n = \frac{1}{T} \int_{-T/2}^{T/2} f(t) e^{-i\omega_n t} \, dt = \frac{\Delta \omega}{2\pi} \int_{-T/2}^{T/2} f(t) e^{-i\omega_n t} \, dt.
\]

As $T \to \infty$, $\Delta \omega \to d\omega$ and the sum becomes an integral.

\subsection{Definition and Basic Properties}

\begin{definition}[Fourier Transform]
For $f \in L^1(\R)$, the \emph{Fourier transform} is
\[
\hat{f}(\omega) = \fourier\{f\}(\omega) = \int_{-\infty}^{\infty} f(t) e^{-i\omega t} \, dt.
\]
The \emph{inverse Fourier transform} is
\[
f(t) = \fourier^{-1}\{\hat{f}\}(t) = \frac{1}{2\pi} \int_{-\infty}^{\infty} \hat{f}(\omega) e^{i\omega t} \, d\omega.
\]
\end{definition}

\begin{remark}
There are several conventions for the Fourier transform differing in the placement of factors of $2\pi$ and the sign of the exponent. We use the angular frequency convention where $\omega = 2\pi f$.
\end{remark}

\begin{theorem}[Riemann-Lebesgue Lemma]
If $f \in L^1(\R)$, then $\hat{f}$ is continuous and
\[
\lim_{\abs{\omega} \to \infty} \hat{f}(\omega) = 0.
\]
\end{theorem}

\begin{proof}
\emph{Continuity.} If $\omega_k \to \omega$, then $f(t)e^{-i\omega_k t} \to f(t)e^{-i\omega t}$ for every $t$, and the integrands are dominated by $\abs{f} \in L^1$; dominated convergence gives $\hat f(\omega_k) \to \hat f(\omega)$.

\emph{Decay.} For the indicator of an interval, $f = \bbone_{[a,b]}$, a direct computation gives $\hat f(\omega) = (e^{-i\omega a} - e^{-i\omega b})/(i\omega)$, so $\abs{\hat f(\omega)} \le 2/\abs{\omega} \to 0$. By linearity the same holds for every step function. Step functions are dense in $L^1(\R)$: given $\varepsilon > 0$ pick a step function $g$ with $\norm{f-g}_1 < \varepsilon$. Since $\abs{\hat f(\omega) - \hat g(\omega)} \le \norm{f-g}_1$ for every $\omega$,
\[
\limsup_{\abs{\omega}\to\infty} \abs{\hat f(\omega)} \le \varepsilon + \lim_{\abs{\omega}\to\infty}\abs{\hat g(\omega)} = \varepsilon,
\]
and $\varepsilon$ is arbitrary.
\end{proof}

\subsection{Inversion Formula}

\begin{theorem}[Fourier Inversion Theorem]
Let $f \in L^1(\R)$ with $\hat{f} \in L^1(\R)$. If $f$ is continuous at $t$, then
\[
f(t) = \frac{1}{2\pi} \int_{-\infty}^{\infty} \hat{f}(\omega) e^{i\omega t} \, d\omega.
\]
\end{theorem}

\begin{remark}
When $\hat f \in L^1$, the right side is a continuous function of $t$ (by the continuity part of the Riemann--Lebesgue argument), so $f$ agrees almost everywhere with a continuous function and cannot have jumps. For a function with jumps, such as $\rect$, $\hat f \notin L^1$; the inversion integral must then be read as the limit $\lim_{R\to\infty}\frac1{2\pi}\int_{-R}^R$, which converges to $(f(t^+)+f(t^-))/2$ at a jump of a piecewise smooth $f$, exactly as for Fourier series.
\end{remark}

\subsection{Plancherel Theorem and $L^2$ Theory}

\begin{theorem}[Plancherel Theorem]\label{fourier-analysis:thm:plancherel}
For $f \in L^1(\R)\cap L^2(\R)$ one has $\hat f \in L^2(\R)$ and the identity below. Consequently the Fourier transform extends by continuity from the dense subspace $L^1\cap L^2$ to a map $L^2(\R) \to L^2(\R)$ which, up to the factor $1/\sqrt{2\pi}$, is unitary:
\[
\int_{-\infty}^{\infty} \abs{f(t)}^2 \, dt = \frac{1}{2\pi} \int_{-\infty}^{\infty} \abs{\hat{f}(\omega)}^2 \, d\omega.
\]
\end{theorem}

\begin{corollary}[Parseval's Identity]
For $f, g \in L^2(\R)$,
\[
\int_{-\infty}^{\infty} f(t) \overline{g(t)} \, dt = \frac{1}{2\pi} \int_{-\infty}^{\infty} \hat{f}(\omega) \overline{\hat{g}(\omega)} \, d\omega.
\]
\end{corollary}

\subsection{Properties of the Fourier Transform}

\begin{proposition}[Linearity]
For $\alpha, \beta \in \C$ and $f, g \in L^1(\R)$,
\[
\fourier\{\alpha f + \beta g\} = \alpha \hat{f} + \beta \hat{g}.
\]
\end{proposition}

\begin{proposition}[Time Shift]
\[
\fourier\{f(t - t_0)\}(\omega) = e^{-i\omega t_0} \hat{f}(\omega).
\]
\end{proposition}

\begin{proof}
\begin{align*}
\fourier\{f(t - t_0)\}(\omega) &= \int_{-\infty}^{\infty} f(t - t_0) e^{-i\omega t} \, dt \\
&= \int_{-\infty}^{\infty} f(u) e^{-i\omega (u + t_0)} \, du \quad (u = t - t_0) \\
&= e^{-i\omega t_0} \int_{-\infty}^{\infty} f(u) e^{-i\omega u} \, du \\
&= e^{-i\omega t_0} \hat{f}(\omega).
\end{align*}
\end{proof}

\begin{proposition}[Frequency Shift / Modulation]
\[
\fourier\{e^{i\omega_0 t} f(t)\}(\omega) = \hat{f}(\omega - \omega_0).
\]
\end{proposition}

\begin{proposition}[Scaling]
For $a \neq 0$,
\[
\fourier\{f(at)\}(\omega) = \frac{1}{\abs{a}} \hat{f}\left(\frac{\omega}{a}\right).
\]
\end{proposition}

\begin{proof}
For $a > 0$:
\begin{align*}
\fourier\{f(at)\}(\omega) &= \int_{-\infty}^{\infty} f(at) e^{-i\omega t} \, dt \\
&= \int_{-\infty}^{\infty} f(u) e^{-i\omega u/a} \frac{du}{a} \quad (u = at) \\
&= \frac{1}{a} \hat{f}\left(\frac{\omega}{a}\right).
\end{align*}
For $a < 0$, the limits reverse and we get a factor of $1/\abs{a}$.
\end{proof}

\begin{proposition}[Differentiation in Time]
If $f$ and $f'$ are in $L^1(\R)$ and $\lim_{\abs{t} \to \infty} f(t) = 0$, then
\[
\fourier\{f'(t)\}(\omega) = i\omega \hat{f}(\omega).
\]
\end{proposition}

\begin{proof}
Using integration by parts:
\begin{align*}
\fourier\{f'(t)\}(\omega) &= \int_{-\infty}^{\infty} f'(t) e^{-i\omega t} \, dt \\
&= \left[ f(t) e^{-i\omega t} \right]_{-\infty}^{\infty} - \int_{-\infty}^{\infty} f(t) (-i\omega) e^{-i\omega t} \, dt \\
&= 0 + i\omega \int_{-\infty}^{\infty} f(t) e^{-i\omega t} \, dt \\
&= i\omega \hat{f}(\omega).
\end{align*}
\end{proof}

\begin{proposition}[Differentiation in Frequency]
If $tf(t) \in L^1(\R)$, then $\hat{f}$ is differentiable and
\[
\frac{d\hat{f}}{d\omega}(\omega) = \fourier\{-it f(t)\}(\omega).
\]
\end{proposition}

\begin{proposition}[Convolution Theorem]
If $f, g \in L^1(\R)$, then
\[
\fourier\{f * g\} = \hat{f} \cdot \hat{g},
\]
where $(f * g)(t) = \int_{-\infty}^{\infty} f(\tau) g(t - \tau) \, d\tau$ is the convolution.
\end{proposition}

\begin{proof}
\begin{align*}
\fourier\{f * g\}(\omega) &= \int_{-\infty}^{\infty} \left( \int_{-\infty}^{\infty} f(\tau) g(t - \tau) \, d\tau \right) e^{-i\omega t} \, dt \\
&= \int_{-\infty}^{\infty} f(\tau) \left( \int_{-\infty}^{\infty} g(t - \tau) e^{-i\omega t} \, dt \right) d\tau \\
&= \int_{-\infty}^{\infty} f(\tau) e^{-i\omega \tau} \left( \int_{-\infty}^{\infty} g(u) e^{-i\omega u} \, du \right) d\tau \\
&= \hat{g}(\omega) \int_{-\infty}^{\infty} f(\tau) e^{-i\omega \tau} \, d\tau \\
&= \hat{f}(\omega) \hat{g}(\omega).
\end{align*}
\end{proof}

\subsection{Important Fourier Transform Pairs}

\begin{example}[Rectangular Pulse]
\[
\rect(t) = \begin{cases}
1 & \abs{t} < 1/2 \\
0 & \abs{t} > 1/2
\end{cases}
\]
Then
\[
\fourier\{\rect(t)\}(\omega) = \int_{-1/2}^{1/2} e^{-i\omega t} \, dt = \frac{2\sin(\omega/2)}{\omega} = \sinc(\omega/2\pi),
\]
where $\sinc(x) = \sin(\pi x)/(\pi x)$.
\end{example}

\begin{example}[Gaussian]
For $f(t) = e^{-\alpha t^2}$ with $\alpha > 0$,
\[
\hat{f}(\omega) = \sqrt{\frac{\pi}{\alpha}} e^{-\omega^2/(4\alpha)}.
\]
The Gaussian is its own Fourier transform (up to scaling).
\end{example}

\begin{proof}
Let $I(\omega) = \int_{-\infty}^{\infty} e^{-\alpha t^2 - i\omega t} \, dt$.

Complete the square:
\[
-\alpha t^2 - i\omega t = -\alpha \left( t + \frac{i\omega}{2\alpha} \right)^2 - \frac{\omega^2}{4\alpha}.
\]

Shift the contour:
\[
I(\omega) = e^{-\omega^2/(4\alpha)} \int_{-\infty}^{\infty} e^{-\alpha u^2} \, du = e^{-\omega^2/(4\alpha)} \sqrt{\frac{\pi}{\alpha}}.
\]
\end{proof}

\begin{example}[Dirac Delta]
The Dirac delta $\delta(t)$ (technically a distribution) has
\[
\hat{\delta}(\omega) = 1.
\]
Conversely,
\[
\fourier\{1\}(\omega) = 2\pi \delta(\omega).
\]
\end{example}

\section{Discrete-Time Fourier Transform (DTFT)}

\subsection{From Continuous to Discrete Time}

When we sample a continuous signal $f(t)$ at intervals $T_s$, we obtain a discrete sequence $x[n] = f(nT_s)$ for $n \in \Z$. The DTFT analyzes such sequences.

\begin{definition}[DTFT]
For a sequence $\{x[n]\}_{n=-\infty}^{\infty}$ with $\sum_{n=-\infty}^{\infty} \abs{x[n]} < \infty$, the \emph{Discrete-Time Fourier Transform} is
\[
X(e^{i\omega}) = \sum_{n=-\infty}^{\infty} x[n] e^{-i\omega n}.
\]
The inverse DTFT is
\[
x[n] = \frac{1}{2\pi} \int_{-\pi}^{\pi} X(e^{i\omega}) e^{i\omega n} \, d\omega.
\]
\end{definition}

\begin{remark}
The frequency variable $\omega$ is now normalized (dimensionless). The actual frequency is $f = \omega/(2\pi T_s)$ Hz.
\end{remark}

\subsection{Periodicity in Frequency}

\begin{proposition}
$X(e^{i\omega})$ is $2\pi$-periodic in $\omega$:
\[
X(e^{i(\omega + 2\pi)}) = X(e^{i\omega}).
\]
\end{proposition}

\begin{proof}
Since $e^{-i2\pi n} = 1$ for every integer $n$,
\[
X(e^{i(\omega + 2\pi)}) = \sum_{n} x[n] e^{-i\omega n} e^{-i 2\pi n} = \sum_{n} x[n] e^{-i\omega n} = X(e^{i\omega}).
\]
\end{proof}

This periodicity reflects the \emph{aliasing} phenomenon: frequencies differing by multiples of the sampling frequency are indistinguishable.

\subsection{Sampling Theorem}

\begin{theorem}[Nyquist-Shannon Sampling Theorem]
Let $f(t)$ be a bandlimited signal with $\hat{f}(\omega) = 0$ for $\abs{\omega} > \omega_{\max}$. If $f$ is sampled at rate $\omega_s = 2\pi/T_s > 2\omega_{\max}$ (the Nyquist rate), then $f$ can be perfectly reconstructed from its samples:
\[
f(t) = \sum_{n=-\infty}^{\infty} f(nT_s) \sinc\left(\frac{t - nT_s}{T_s}\right).
\]
\end{theorem}

Here and below $\sinc(x) = \sin(\pi x)/(\pi x)$ (with $\sinc(0) = 1$) is the \emph{normalised} sinc, which vanishes at every nonzero integer. The interpolating function $\sinc((t-nT_s)/T_s)$ is therefore $1$ at $t = nT_s$ and $0$ at all other sample times, so the formula reproduces the samples exactly. A proof is given in the chapter appendix.

\subsection{Relationship to Fourier Transform}

If $x[n] = f(nT_s)$, where $f$ is regular enough for the Poisson summation formula to hold (for instance $f$ and $\hat f$ both decaying at least like $(1+\abs{t})^{-2}$), then
\[
X(e^{i\omega}) = \frac{1}{T_s} \sum_{k=-\infty}^{\infty} \hat{f}\left(\frac{\omega - 2\pi k}{T_s}\right).
\]

This shows that the DTFT is a periodically replicated version of the continuous Fourier transform. If $\hat f$ vanishes outside $\abs{\omega} < \pi/T_s$, the replicas do not overlap; otherwise they add up, and high frequencies masquerade as low ones (aliasing). This is the content of the sampling theorem.

\subsection{Properties of the DTFT}

\begin{proposition}[Time Shift]
\[
\text{DTFT}\{x[n - n_0]\} = e^{-i\omega n_0} X(e^{i\omega}).
\]
\end{proposition}

\begin{proposition}[Frequency Shift]
\[
\text{DTFT}\{e^{i\omega_0 n} x[n]\} = X(e^{i(\omega - \omega_0)}).
\]
\end{proposition}

\begin{proposition}[Convolution]
For discrete convolution $(x * h)[n] = \sum_{k=-\infty}^{\infty} x[k] h[n-k]$,
\[
\text{DTFT}\{x * h\} = X(e^{i\omega}) H(e^{i\omega}).
\]
\end{proposition}

\begin{proposition}[Parseval's Identity]
\[
\sum_{n=-\infty}^{\infty} \abs{x[n]}^2 = \frac{1}{2\pi} \int_{-\pi}^{\pi} \abs{X(e^{i\omega})}^2 \, d\omega.
\]
\end{proposition}

\section{Discrete Fourier Transform (DFT)}

\subsection{Finite-Length Sequences}

In practice, we work with finite-length sequences. The DFT is the appropriate tool for this setting.

\begin{definition}[DFT]\label{fourier-analysis:def:dft}
For a finite sequence $\{x[n]\}_{n=0}^{N-1}$ of length $N$, the \emph{Discrete Fourier Transform} is
\[
X[k] = \sum_{n=0}^{N-1} x[n] e^{-i 2\pi kn/N}, \quad k = 0, 1, \ldots, N-1.
\]
The \emph{inverse DFT} (IDFT) is
\[
x[n] = \frac{1}{N} \sum_{k=0}^{N-1} X[k] e^{i 2\pi kn/N}, \quad n = 0, 1, \ldots, N-1.
\]
\end{definition}

\subsection{Matrix Formulation}

Define the \emph{DFT matrix}:
\[
W_N = e^{-i 2\pi/N}, \quad [F_N]_{kn} = W_N^{kn} = e^{-i 2\pi kn/N}.
\]

Then:
\[
\mathbf{X} = F_N \mathbf{x}, \quad \mathbf{x} = \frac{1}{N} F_N^* \mathbf{X},
\]
where $F_N^*$ is the conjugate transpose.

\begin{lemma}
$F_N$ is unitary up to scaling: $F_N^* F_N = N I$.
\end{lemma}

\begin{proof}
\begin{align*}
[F_N^* F_N]_{mn} &= \sum_{k=0}^{N-1} \overline{W_N^{km}} W_N^{kn} = \sum_{k=0}^{N-1} W_N^{k(n-m)} \\
&= \begin{cases}
N & \text{if } m = n \\
\frac{1 - W_N^{N(n-m)}}{1 - W_N^{n-m}} = \frac{1 - 1}{1 - W_N^{n-m}} = 0 & \text{if } m \neq n.
\end{cases}
\end{align*}
\end{proof}

\subsection{Relationship to DTFT and Fourier Series}

\begin{proposition}
The DFT is the DTFT sampled at $N$ equally-spaced frequencies:
\[
X[k] = X(e^{i\omega})\big|_{\omega = 2\pi k/N}.
\]
\end{proposition}

Alternatively, if we periodically extend $x[n]$ to create $\tilde{x}[n] = x[n \mod N]$, then $X[k]$ are the Fourier series coefficients of $\tilde{x}[n]$ (up to scaling).

\subsection{Circular Convolution}

\begin{definition}[Circular Convolution]\label{fourier-analysis:def:circular-convolution}
The \emph{circular convolution} of two length-$N$ sequences is
\[
(x \circledast h)[n] = \sum_{m=0}^{N-1} x[m] h[(n-m) \mod N].
\]
\end{definition}

\begin{theorem}[Circular Convolution Theorem]\label{fourier-analysis:thm:circular-convolution}
\[
\text{DFT}\{x \circledast h\} = \text{DFT}\{x\} \cdot \text{DFT}\{h\}.
\]
\end{theorem}

\begin{remark}
To perform linear convolution using the DFT, we must zero-pad the sequences to length at least $N_1 + N_2 - 1$ to avoid time-domain aliasing.
\end{remark}

\subsection{Properties of the DFT}

\begin{proposition}[Parseval's Identity]
\[
\sum_{n=0}^{N-1} \abs{x[n]}^2 = \frac{1}{N} \sum_{k=0}^{N-1} \abs{X[k]}^2.
\]
\end{proposition}

\begin{proposition}[Symmetry for Real Sequences]
If $x[n] \in \R$, then $X[N-k] = \overline{X[k]}$.
\end{proposition}

\section{Fast Fourier Transform (FFT)}

\subsection{Computational Complexity}

Computing the DFT directly requires $O(N^2)$ complex multiplications. The FFT is a class of algorithms that compute the DFT in $O(N \log N)$ operations.

\subsection{Cooley-Tukey FFT Algorithm}

Assume $N = 2^m$ (radix-2 FFT). We use the \emph{decimation-in-time} approach.

\subsubsection{Derivation}

Split the sum into even and odd indices:
\begin{align*}
X[k] &= \sum_{n=0}^{N-1} x[n] W_N^{kn} \\
&= \sum_{r=0}^{N/2-1} x[2r] W_N^{k(2r)} + \sum_{r=0}^{N/2-1} x[2r+1] W_N^{k(2r+1)} \\
&= \sum_{r=0}^{N/2-1} x[2r] W_{N/2}^{kr} + W_N^k \sum_{r=0}^{N/2-1} x[2r+1] W_{N/2}^{kr} \\
&= E[k] + W_N^k O[k],
\end{align*}
where $E[k]$ and $O[k]$ are the $(N/2)$-point DFTs of the even and odd indexed samples.

\begin{lemma}
$W_N^{k+N/2} = -W_N^k$ and $E[k + N/2] = E[k]$, $O[k + N/2] = O[k]$.
\end{lemma}

Therefore:
\begin{align*}
X[k] &= E[k] + W_N^k O[k], \quad k = 0, \ldots, N/2-1, \\
X[k + N/2] &= E[k] - W_N^k O[k], \quad k = 0, \ldots, N/2-1.
\end{align*}

\subsubsection{Butterfly Computation}

The basic operation, called a \emph{butterfly}, is:
\begin{align*}
A &= E + W \cdot O, \\
B &= E - W \cdot O.
\end{align*}

This requires one complex multiplication and two complex additions.

\subsubsection{Computational Cost}

Let $C(N)$ be the number of complex additions and multiplications used for an $N$-point DFT by this recursion. The two half-size DFTs cost $2C(N/2)$; combining them takes $N/2$ butterflies, each with one multiplication and two additions, so at most $cN$ further operations for a constant $c$ ($c = 3/2$ here). With $C(1) = 0$,
\[
C(N) = 2C(N/2) + cN \quad\Longrightarrow\quad C(N) = cN\log_2 N,
\]
as one checks by induction: $2c\frac N2\log_2\frac N2 + cN = cN(\log_2 N - 1) + cN = cN\log_2 N$. In particular the radix-2 FFT uses $\frac N2\log_2 N$ complex multiplications, against $N^2$ for the direct sum; for $N = 1024$ that is $5120$ instead of about $10^6$.

\subsection{Bit-Reversal Permutation}

If the recursion is unrolled into an in-place iterative algorithm, the input must be reordered according to bit-reversal: after $\log_2 N$ even/odd splits, the sample $x[n]$ ends up in the position whose binary expansion is that of $n$ read backwards.

\begin{example}
For $N = 8$:
\begin{center}
\begin{tabular}{c|c|c}
Index & Binary & Bit-Reversed \\
\hline
0 & 000 & 000 = 0 \\
1 & 001 & 100 = 4 \\
2 & 010 & 010 = 2 \\
3 & 011 & 110 = 6 \\
4 & 100 & 001 = 1 \\
5 & 101 & 101 = 5 \\
6 & 110 & 011 = 3 \\
7 & 111 & 111 = 7 \\
\end{tabular}
\end{center}
\end{example}

\subsection{Variants and Generalizations}

\begin{itemize}
\item \textbf{Radix-4, Radix-8, etc.}: Further subdivide into 4 or 8 parts for better cache performance.
\item \textbf{Split-radix FFT}: Combines radix-2 and radix-4 for optimal operation count.
\item \textbf{Mixed-radix FFT}: For $N$ that is not a power of 2.
\item \textbf{Prime-factor algorithm}: For $N = N_1 N_2$ with $\gcd(N_1, N_2) = 1$.
\item \textbf{Bluestein's algorithm}: For arbitrary $N$ using chirp z-transform.
\end{itemize}

\section{The four transforms side by side}


\begin{center}
\small
\resizebox{\linewidth}{!}{%
\begin{tabular}{|l|c|c|c|c|}
\hline
\textbf{Transform} & \textbf{Time} & \textbf{Frequency} & \textbf{Forward} & \textbf{Inverse} \\
\hline
Fourier & Cont., & Discrete, & $c_n = \frac{1}{T}\int_0^T f(t) e^{-i2\pi nt/T} dt$ & $f(t) = \sum_n c_n e^{i2\pi nt/T}$ \\
Series (FS) & Periodic & Aperiodic & & \\
\hline
Fourier & Cont., & Cont., & $\hat{f}(\omega) = \int f(t) e^{-i\omega t} dt$ & $f(t) = \frac{1}{2\pi}\int \hat{f}(\omega) e^{i\omega t} d\omega$ \\
Transform (FT) & Aperiodic & Aperiodic & & \\
\hline
DTFT & Discrete, & Cont., & $X(e^{i\omega}) = \sum_n x[n] e^{-i\omega n}$ & $x[n] = \frac{1}{2\pi}\int_{-\pi}^{\pi} X(e^{i\omega}) e^{i\omega n} d\omega$ \\
& Aperiodic & Periodic & & \\
\hline
DFT & Discrete, & Discrete, & $X[k] = \sum_{n=0}^{N-1} x[n] e^{-i2\pi kn/N}$ & $x[n] = \frac{1}{N}\sum_{k=0}^{N-1} X[k] e^{i2\pi kn/N}$ \\
& Periodic & Periodic & & \\
\hline
FFT & \multicolumn{4}{c|}{Efficient algorithm for computing DFT in $O(N \log N)$ time} \\
\hline
\end{tabular}%
}
\end{center}

\paragraph{When each transform exists.}

\begin{center}
\begin{tabular}{|l|p{0.75\linewidth}|}
\hline
\textbf{FS} & $f \in L^2([0,T])$ for $L^2$ convergence; Dirichlet conditions for pointwise convergence \\
\hline
\textbf{FT} & $f \in L^1(\R)$ for existence; $f \in L^2(\R)$ for Plancherel; additional smoothness for pointwise convergence \\
\hline
\textbf{DTFT} & $\sum_n \abs{x[n]} < \infty$ (absolutely summable sequences) \\
\hline
\textbf{DFT} & Always exists for finite sequences \\
\hline
\end{tabular}
\end{center}

\section{Applications and Examples}

\subsection{Solving Differential Equations}

\begin{example}[Heat Equation]
Consider the heat equation on $\R$:
\[
\frac{\partial u}{\partial t} = \alpha \frac{\partial^2 u}{\partial x^2}, \quad u(x, 0) = f(x).
\]

Apply the Fourier transform in $x$:
\[
\frac{\partial \hat{u}}{\partial t} = -\alpha \omega^2 \hat{u}.
\]

This is an ODE with solution:
\[
\hat{u}(\omega, t) = \hat{f}(\omega) e^{-\alpha \omega^2 t}.
\]

Inverse transform:
\[
u(x, t) = \frac{1}{2\pi} \int \hat{f}(\omega) e^{-\alpha \omega^2 t} e^{i\omega x} d\omega = (f * G_t)(x),
\]
where $G_t(x) = \frac{1}{\sqrt{4\pi \alpha t}} e^{-x^2/(4\alpha t)}$ is the heat kernel.
\end{example}

\subsection{Signal Processing}

\begin{example}[Ideal Low-Pass Filter]
An ideal low-pass filter has frequency response:
\[
H(\omega) = \begin{cases}
1 & \abs{\omega} < \omega_c \\
0 & \abs{\omega} > \omega_c
\end{cases}
\]

The impulse response is:
\[
h(t) = \frac{1}{2\pi} \int_{-\omega_c}^{\omega_c} e^{i\omega t} d\omega = \frac{\omega_c}{\pi} \sinc(\omega_c t/\pi).
\]

Note that $h(t)$ is non-causal and extends to $\pm \infty$, making it unrealizable in practice.
\end{example}

\subsection{Fast Convolution}

For long sequences, computing convolution directly takes $O(N^2)$ time. Using the FFT:
\begin{enumerate}
\item Zero-pad $x$ and $h$ to length $N \geq N_x + N_h - 1$.
\item Compute $X = \text{FFT}(x)$ and $H = \text{FFT}(h)$.
\item Compute $Y = X \cdot H$ (elementwise).
\item Compute $y = \text{IFFT}(Y)$.
\end{enumerate}

Total complexity: $O(N \log N)$.

\section{Advanced Topics}

\subsection{Windowing and Spectral Leakage}

When applying the DFT to a finite segment of a signal, we implicitly multiply by a rectangular window, causing spectral leakage.

\begin{definition}[Window Functions]
Common windows include:
\begin{itemize}
\item \textbf{Rectangular}: $w[n] = 1$
\item \textbf{Hamming}: $w[n] = 0.54 - 0.46\cos(2\pi n/(N-1))$
\item \textbf{Hann}: $w[n] = 0.5(1 - \cos(2\pi n/(N-1)))$
\item \textbf{Blackman}: $w[n] = 0.42 - 0.5\cos(2\pi n/(N-1)) + 0.08\cos(4\pi n/(N-1))$
\end{itemize}
\end{definition}

Windows trade off main-lobe width (frequency resolution) against side-lobe level (leakage).

\subsection{Short-Time Fourier Transform (STFT)}

For non-stationary signals, the STFT provides time-frequency localization:
\[
\text{STFT}_x(\tau, \omega) = \int x(t) w(t - \tau) e^{-i\omega t} dt,
\]
where $w(t)$ is a window function.

\subsection{Discrete Cosine Transform (DCT)}

The DCT is widely used in image and audio compression (JPEG, MP3):
\[
X[k] = \sum_{n=0}^{N-1} x[n] \cos\left(\frac{\pi k(2n+1)}{2N}\right).
\]

The DCT is real-valued and has excellent energy compaction properties for natural signals.

\subsection{Multidimensional Transforms}

The DFT extends naturally to higher dimensions:
\[
X[k_1, k_2] = \sum_{n_1=0}^{N_1-1} \sum_{n_2=0}^{N_2-1} x[n_1, n_2] e^{-i2\pi(k_1 n_1/N_1 + k_2 n_2/N_2)}.
\]

The 2D FFT is computed by applying 1D FFTs along each dimension (row-column algorithm).


\chapterappendix

\section{Proofs of Advanced Theorems}

\subsection{Proof of Plancherel Theorem}

\begin{proof}[Proof of Theorem~\ref{fourier-analysis:thm:plancherel}]
Let $f \in L^1(\R) \cap L^2(\R)$; the extension to $L^2(\R)$ is then a density argument. Define $\tilde f(t) = \overline{f(-t)}$ and the autocorrelation
\[
g(t) = (f * \tilde f)(t) = \int f(s)\,\overline{f(s-t)}\,ds .
\]
Then $g \in L^1$ (convolution of two $L^1$ functions), $g$ is bounded and continuous (by Cauchy--Schwarz and continuity of translation in $L^2$), and $g(0) = \norm{f}_2^2$. By the convolution theorem and $\hat{\tilde f} = \overline{\hat f}$, $\hat g = \hat f\,\overline{\hat f} = \abs{\hat f}^2 \ge 0$.

Fix $\varepsilon > 0$. By Fubini and the Gaussian pair (Example on the Gaussian, with $\alpha = 1/(4\varepsilon)$),
\[
\frac{1}{2\pi}\int \hat g(\omega)\,e^{-\varepsilon\omega^2}d\omega
= \int g(t)\,G_\varepsilon(t)\,dt,
\qquad G_\varepsilon(t) = \frac{1}{\sqrt{4\pi\varepsilon}}e^{-t^2/(4\varepsilon)} .
\]
The kernels $G_\varepsilon$ are nonnegative, integrate to $1$ and concentrate at $0$ as $\varepsilon\to0$; since $g$ is bounded and continuous at $0$, the right side tends to $g(0) = \norm{f}_2^2$. On the left, $\abs{\hat f}^2 e^{-\varepsilon\omega^2}$ increases to $\abs{\hat f}^2$, so by monotone convergence the left side tends to $\frac1{2\pi}\int\abs{\hat f}^2$. Hence $\norm{f}_2^2 = \frac1{2\pi}\norm{\hat f}_2^2$, which in particular shows $\hat f \in L^2$.
\end{proof}

\begin{remark}
A tempting shortcut is to apply the inversion formula to the truncations $f\bbone_{[-R,R]}$ and let $R\to\infty$. That argument is circular: inversion requires $\hat f \in L^1$, which is not known for a general $f \in L^1\cap L^2$. The Gaussian factor $e^{-\varepsilon\omega^2}$ above is what makes every integral absolutely convergent.
\end{remark}

\subsection{Proof of Shannon Sampling Theorem}

\begin{proof}
Since $f$ is bandlimited, $\hat{f}(\omega) = 0$ for $\abs{\omega} > \omega_{\max}$. Extend $\hat{f}$ periodically with period $\omega_s > 2\omega_{\max}$:
\[
\hat{f}_p(\omega) = \sum_{k=-\infty}^{\infty} \hat{f}(\omega - k\omega_s).
\]

Since $\omega_s > 2\omega_{\max}$, there is no overlap and $\hat{f}_p(\omega) = \hat{f}(\omega)$ for $\abs{\omega} < \omega_s/2$.

The Fourier series coefficients of $\hat{f}_p$ are:
\[
c_n = \frac{1}{\omega_s} \int_{-\omega_s/2}^{\omega_s/2} \hat{f}(\omega) e^{i\omega n T_s} d\omega = \frac{1}{\omega_s} \int \hat{f}(\omega) e^{i\omega n T_s} d\omega = \frac{2\pi}{\omega_s} f(nT_s).
\]

Therefore:
\[
\hat{f}(\omega) = \hat{f}_p(\omega) = \sum_n c_n e^{-i\omega n T_s} = \frac{2\pi}{\omega_s} \sum_n f(nT_s) e^{-i\omega n T_s}.
\]

Inverse Fourier transform (the factor $\frac{1}{2\pi}\cdot\frac{2\pi}{\omega_s} = \frac1{\omega_s}$):
\[
f(t) = \frac{1}{2\pi} \int \hat{f}(\omega) e^{i\omega t} d\omega = \sum_n f(nT_s) \frac{1}{\omega_s} \int_{-\omega_s/2}^{\omega_s/2} e^{i\omega(t - nT_s)} d\omega.
\]

The integral evaluates to:
\[
\frac{1}{\omega_s} \int_{-\omega_s/2}^{\omega_s/2} e^{i\omega(t - nT_s)} d\omega = \frac{\sin(\omega_s(t - nT_s)/2)}{\omega_s(t - nT_s)/2} = \sinc\left(\frac{t - nT_s}{T_s}\right),
\]
using $\omega_s/2 = \pi/T_s$. Substituting gives the interpolation formula. (Exchanging sum and integral is justified in $L^2([-\omega_s/2,\omega_s/2])$, where the Fourier series of $\hat f$ converges.)
\end{proof}

\begin{chaptersummary}
\begin{itemize}
\item The Fourier series of a $T$-periodic $f \in L^2([0,T])$ has coefficients $c_n = \frac1T\int_0^T f(t)e^{-i2\pi nt/T}dt$, obtained from the orthogonality of the exponentials; it converges in mean square, pointwise to the midpoint of jumps under the Dirichlet conditions, and Parseval's identity $\frac1T\int_0^T\abs{f}^2 = \sum_n\abs{c_n}^2$ holds.
\item The Fourier transform $\hat f(\omega) = \int f(t)e^{-i\omega t}dt$ is the $T\to\infty$ limit of the series. For $f\in L^1$ it is continuous and vanishes at infinity (Riemann--Lebesgue); Plancherel extends it to $L^2$, where $\norm{f}_2^2 = \frac1{2\pi}\norm{\hat f}_2^2$.
\item Shifts become phase factors, derivatives become multiplication by $i\omega$, and convolution becomes multiplication: $\widehat{f*g} = \hat f\hat g$.
\item Sampling at interval $T_s$ replicates the spectrum with period $2\pi/T_s$ (DTFT); a signal bandlimited below the Nyquist frequency $\pi/T_s$ is recovered exactly by sinc interpolation, otherwise aliasing occurs.
\item The DFT is the matrix $F_N$ with $F_N^*F_N = NI$; it samples the DTFT at $N$ frequencies, diagonalises circular convolution, and is computed by the FFT in $O(N\log N)$ operations by recursive even/odd splitting.
\end{itemize}
\end{chaptersummary}

\section*{Further reading}
\addcontentsline{toc}{section}{Further reading}
Stein and Shakarchi (\emph{Fourier Analysis: An Introduction}, 2003) develop Fourier series and the transform with complete proofs at the level of this chapter; Rudin (\emph{Real and Complex Analysis}, 3rd ed., 1987) gives the $L^1$/$L^2$ theory, including Plancherel. Oppenheim and Schafer (\emph{Discrete-Time Signal Processing}, 3rd ed., 2009) is the standard reference for the DTFT, the DFT, sampling and windowing, and Bracewell (\emph{The Fourier Transform and Its Applications}, 3rd ed., 2000) is a rich source of transform pairs and intuition. The FFT in its modern form is due to Cooley and Tukey (\emph{An Algorithm for the Machine Calculation of Complex Fourier Series}, 1965).

\section*{Exercises}
\addcontentsline{toc}{section}{Exercises}

\begin{exercise}\label{fourier-analysis:exr:1}
Compute the complex Fourier coefficients of the sawtooth $f(t) = t$ on $(-\pi,\pi)$, extended $2\pi$-periodically, write the real form of the series, and use Parseval's identity to evaluate $\sum_{n\ge1} 1/n^2$.
\hint{Integrate by parts; $c_0 = 0$ because $f$ is odd.}
\solutionlink{fourier-analysis:exr:1}
\end{exercise}

\begin{exercise}\label{fourier-analysis:exr:2}
Show that for $a > 0$ the Fourier transform of $f(t) = e^{-a\abs{t}}$ is $\hat f(\omega) = 2a/(a^2+\omega^2)$. Then use Plancherel to compute $\int_{-\infty}^{\infty} (a^2+\omega^2)^{-2}\,d\omega$.
\solutionlink{fourier-analysis:exr:2}
\end{exercise}

\begin{exercise}\label{fourier-analysis:exr:3}
Prove that if $f\in L^1(\R)$ is real and even then $\hat f$ is real and even, and if $f$ is real and odd then $\hat f$ is purely imaginary and odd.
\solutionlink{fourier-analysis:exr:3}
\end{exercise}

\begin{exercise}\label{fourier-analysis:exr:4}
Let $x = (1, 2, 0, 0)$ and $h = (1, 1, 0, 0)$, viewed as length-$4$ sequences.
\begin{enumerate}
\item Compute $X = F_4x$ and $H = F_4h$.
\item Compute $y = x \circledast h$ directly, and check that $F_4 y = X\cdot H$ entrywise.
\item Compute $x \circledast h$ for the unpadded length-$2$ sequences $(1,2)$ and $(1,1)$ and explain why it differs from the linear convolution $(1,3,2)$.
\end{enumerate}
\hint{$W_4 = e^{-i\pi/2} = -i$, so the rows of $F_4$ are $(1,1,1,1)$, $(1,-i,-1,i)$, $(1,-1,1,-1)$, $(1,i,-1,-i)$.}
\solutionlink{fourier-analysis:exr:4}
\end{exercise}

\begin{exercise}\label{fourier-analysis:exr:5}
A continuous signal $f(t) = \cos(2\pi\cdot 7t)$ (frequency $7$\,Hz) is sampled at $10$ samples per second. Show that the samples coincide with those of a cosine of a lower frequency, identify that frequency, and state the smallest sampling rate that avoids this.
\solutionlink{fourier-analysis:exr:5}
\end{exercise}

\end{notechapter}
